\section{ANUGA Baseline}
\label{sec:anuga-baseline}

\subsection{Domain construction and parameter assignment}

ANUGA was used as the unstructured-grid numerical reference, rather than as
observational ground truth.  The implementation solves the depth-integrated
shallow-water equations by a finite-volume method on triangular cells
\cite{MungkasiRoberts2012}.  The reprojected mesh contained 82,921 triangular
elements and 41,888 unique nodes in EPSG:32643.  Triangle orientation and
non-degeneracy were checked before the connectivity was passed to
\texttt{anuga.Domain}.

The synthetic DEM was bilinearly sampled at triangle centroids, with
nearest-cell fallback only where bilinear interpolation was undefined.  The
assigned centroid elevations ranged from 216.6333 to 233.4016~m.  Initial
stage was set equal to bed elevation, corresponding to the configured initial
depth of 0.0~m and zero momentum.

Manning resistance was assigned by spatially joining each centroid to the LULC
layer.  The documented effective values were 0.030 for Water, 0.100 for
Built-up, 0.025 for Bare/Open, 0.050 for Low vegetation, and 0.120 for Tree
cover.  These values are traceable to the ranges compiled by
Chow~\cite{Chow1959}, but remain LULC-scale assumptions rather than calibrated
site coefficients.  The source \texttt{roughness} field was required to match
the lookup before assignment.  The element counts were 1,059 Water, 30,606
Built-up, 1,632 Bare/Open, 38,401 Low vegetation, and 11,223 Tree-cover
elements.

\subsection{Numerical and forcing configuration}

The final run used ANUGA's \texttt{DE1} flow algorithm, a CFL number of 0.9,
a 10.0~s yield interval, a 60.0~s output interval, and a minimum storable
height of 0.001~m.  Simulation duration was 3600.0~s.  Rainfall was applied
uniformly to all domain triangles through a rate operator.  The cleaned input
was checked to have units of mm~hr\(^{-1}\) and converted to m~s\(^{-1}\)
with the exact factor \(10^{-3}/3600\).

The source rainfall record contained 35 hourly samples over 122,400~s.  The
one-hour simulation used only 2.9412\% of that record and integrated
2.78499~mm of rainfall.  It did not include the record maximum of
12.141~mm~hr\(^{-1}\), which occurs at 122,400~s.  Accordingly, the run is a
development-window experiment and not a simulation of the complete
\texttt{45rp\_rain} record.

\subsection{Final boundary policy}

All supplied boundary segments and all unmatched exterior mesh edges were
assigned ANUGA reflective boundaries.  Segment tagging used a 1.0~m midpoint
matching tolerance and associated 6, 65, and 10 exterior edges with
\texttt{2Dboundary\_1}, \texttt{2Dboundary\_3}, and
\texttt{2Dboundary\_4}, respectively.  A further 772 exterior edges retained
the generic exterior tag.  Every tag received the same no-normal-flow
condition under the finalized policy
\texttt{all\_reflective\_final}.

\subsection{Boundary Condition Limitations}

The raw attribute schema could not be authoritatively resolved.  The three
records contained
\((\texttt{BndTypeNo},\texttt{ConstValue})=(2,216)\) for
\texttt{2Dboundary\_1} and \((4,0)\) for each of
\texttt{2Dboundary\_3} and \texttt{2Dboundary\_4}.  Inspection of all 47 DBF
fields, available sidecars, repository history, geometry, and elevations
provided no legend or numeric-code mapping.  Field names suggested a possible
DHI MIKE-family export, and the value 216 could plausibly be associated with a
constant-level definition, but neither hypothesis was supported by an
available project database or authoritative schema.  It was therefore not
used to infer boundary behaviour.

The reflective treatment is a conservative, documented simplification rather
than a recovered source-data meaning.  In particular, the short
\texttt{2Dboundary\_1} segment lies on the western to north-western perimeter,
has a synthetic-DEM range of 221.975--222.422~m, and carries a distinct
\texttt{ConstValue}.  If that segment was originally intended to permit
drainage, closing it necessarily underestimates boundary outflow and retains
water that would otherwise leave the domain; this could overestimate local
storage or ponding.  The direction and magnitude of the real hydraulic effect
cannot be quantified without the missing schema and boundary observations.
Internal consistency is maintained because ANUGA, JAX, and the hybrid rollout
all use the same closed policy.

\subsection{Output and run diagnostics}

The finalized run wrote
\texttt{run\_20260727T100639155109Z.sww} and a per-yield mass-balance log
containing 361 records.  Wall-clock runtime was 57.0764~s.  The maximum
absolute mass-balance error was
\(7.23468\times10^{-5}\)~m\(^3\), the maximum relative error was
\(8.52467\times10^{-10}\), and the final error was
\(-5.09780\times10^{-5}\)~m\(^3\).  The maximum absolute recorded boundary
flux was \(1.71811\times10^{-11}\)~m\(^3\), consistent with the reflective
configuration to numerical precision.

SWW output was barycentrically interpolated to the common 50.0~m comparison
grid and saved as \texttt{data/processed/anuga\_baseline.nc}.  The resulting
517 by 551 grid contained 172,841 finite in-mesh cells and 61 output times.
Maximum depth and velocity in this processed development run were
0.0159302~m and 0.106823~m~s\(^{-1}\), respectively.  No non-finite values
were present within the processed mesh mask.
